\documentclass[10pt,a4paper]{amsart}
\usepackage{amsmath,amssymb,amsthm,graphicx}
% For commutative diagrams we use Xy-pic
\usepackage[all]{xy}

% We use 2cell for 2-commutative diagrams.
\xyoption{2cell}
% Theorem environments.
%
\theoremstyle{plain}
\newtheorem{theorem}[subsection]{Theorem}
\newtheorem{proposition}[subsection]{Proposition}
\newtheorem{lemma}[subsection]{Lemma}

\theoremstyle{definition}
\newtheorem{definition}[subsection]{Definition}
\newtheorem{example}[subsection]{Example}
\newtheorem{exercise}[subsection]{Exercise}
\newtheorem{situation}[subsection]{Situation}

\theoremstyle{remark}
\newtheorem{remark}[subsection]{Remark}
\newtheorem{remarks}[subsection]{Remarks}

\numberwithin{equation}{subsection}

% Macros
%
\def\lim{\mathop{\mathrm{lim}}\nolimits}
\def\colim{\mathop{\mathrm{colim}}\nolimits}
\def\Spec{\mathop{\mathrm{Spec}}}
\def\Hom{\mathop{\mathrm{Hom}}\nolimits}
\def\Ext{\mathop{\mathrm{Ext}}\nolimits}
\def\SheafHom{\mathop{\mathcal{H}\!\mathit{om}}\nolimits}
\def\SheafExt{\mathop{\mathcal{E}\!\mathit{xt}}\nolimits}
\def\Sch{\mathit{Sch}}
\def\Mor{\mathop{\mathrm{Mor}}\nolimits}
\def\Ob{\mathop{\mathrm{Ob}}\nolimits}
\def\Sh{\mathop{\mathit{Sh}}\nolimits}
\def\NL{\mathop{N\!L}\nolimits}
\def\CH{\mathop{\mathrm{CH}}\nolimits}
\def\proetale{{pro\text{-}\acute{e}tale}}
\def\etale{{\acute{e}tale}}
\def\QCoh{\mathit{QCoh}}
\def\Ker{\mathop{\mathrm{Ker}}}
\def\Im{\mathop{\mathrm{Im}}}
\def\Coker{\mathop{\mathrm{Coker}}}
\def\Coim{\mathop{\mathrm{Coim}}}

% Boxtimes
%
\DeclareMathSymbol{\boxtimes}{\mathbin}{AMSa}{"02}

%
% Macros for moduli stacks/spaces
%
\def\QCohstack{\mathcal{QC}\!\mathit{oh}}
\def\Cohstack{\mathcal{C}\!\mathit{oh}}
\def\Spacesstack{\mathcal{S}\!\mathit{paces}}
\def\Quotfunctor{\mathrm{Quot}}
\def\Hilbfunctor{\mathrm{Hilb}}
\def\Curvesstack{\mathcal{C}\!\mathit{urves}}
\def\Polarizedstack{\mathcal{P}\!\mathit{olarized}}
\def\Complexesstack{\mathcal{C}\!\mathit{omplexes}}
% \Pic is the operator that assigns to X its picard group, usage \Pic(X)
% \Picardstack_{X/B} denotes the Picard stack of X over B
% \Picardfunctor_{X/B} denotes the Picard functor of X over B
\def\Pic{\mathop{\mathrm{Pic}}\nolimits}
\def\Picardstack{\mathcal{P}\!\mathit{ic}}
\def\Picardfunctor{\mathrm{Pic}}
\def\Deformationcategory{\mathcal{D}\!\mathit{ef}}
\title{The Brauer group of a C1 field is zero}
\author{The Stacks Project Authors}
\date{stacks-project:etale-cohomology.tex:local-snapshot-20261011; GNU FDL 1.2 or later}
\begin{document}
\maketitle
\noindent\textit{Editorial note.} All delivered mathematical text is the author's own, copied from the pinned Stacks Project chapter file etale-cohomology.tex; no statement or proof step was rewritten, reordered or removed.\par\medskip
\noindent\textit{Editorial note.} The proof is the author's complete argument. Its inputs are the author's earlier results of the same section, stated in the proof as 'We have seen that'; the counted span contains no label reference and no bibliography citation, so no navigation stub was needed.\par\medskip
\noindent\textit{Editorial note.} The label theorem-C1-brauer-group-zero occurs in no other staged Stacks chapter, so no cross-chapter restatement had to be weighed for this card.\par\medskip
\noindent\textit{Editorial note.} The preamble is the chapter's own native preamble excerpt already used for this source: the Xy-pic setup (native preamble lines 14-18) and the native theorem declarations with the mathematical macros (native preamble lines 154-217).\par\medskip
\noindent\textit{Allowed standard prerequisites (editorial).}\begin{itemize}
\item C1 field (quasi-algebraically closed field). Same chapter, the definition of a C\_1 field: every homogeneous form of degree d in more than d variables has a nontrivial zero. Supplies the zero of the reduced norm that the proof produces when the degree d exceeds 1.
\item Central simple algebras over a field and their Brauer group. Same chapter, the identification of Br(K) with the set of classes of finite dimensional central division algebras over K. Reduces Br(K) = 0 to the statement that no finite dimensional division algebra with centre K of degree greater than one exists.
\item Splitting of a central division algebra by the separable closure and the reduced norm. Same chapter, the fact that D tensor\_K K\textasciicircum{}sep is isomorphic to Mat\_d(K\textasciicircum{}sep) uniquely up to inner isomorphism, so that the determinant descends. Produces the Galois invariant determinant, hence the homogeneous degree d reduced norm map N\_red : D -> K used in the counted argument.
\end{itemize}\medskip
\begin{theorem}
\label{theorem-C1-brauer-group-zero}
Let $K$ be a $C_1$ field. Then $\text{Br}(K) = 0$.
\end{theorem}

\begin{proof}
Let $D$ be a finite dimensional division algebra over $K$ with center $K$. We
have seen that
$$
D \otimes_K K^{sep} \cong \text{Mat}_d(K^{sep})
$$
uniquely up to inner isomorphism. Hence the determinant $\det :
\text{Mat}_d(K^{sep}) \to K^{sep}$ is Galois invariant and descends to a
homogeneous degree $d$ map
$$
\det = N_\text{red} : D \longrightarrow K
$$
called the {\it reduced norm}. Since $K$ is $C_1$, if $d > 1$, then there
exists a nonzero $x \in D$ with $N_\text{red}(x) = 0$. This clearly implies
that $x$ is not invertible, which is a contradiction. Hence $\text{Br}(K) = 0$.
\end{proof}

\end{document}
